%% kapa: exam=1 round=32 year=2021 period=1 form=A subject=civil_law subject_ko=민법 seq=5 number=5
%% source: Q-Net 감정평가사 기출문제 (cst00309) artlSeq=5208397 file=extracted/1008/5208397/A형/2021년도 제32회 감정평가사 1차 1교시 A형.pdf
%% conversion: kordoc 4.24.51 -> tools/gen1.py (auto); review=proofread-sonnet; flags=none
\begin{question}{5}
법인에 관한 설명으로 옳지 않은 것은? (다툼이 있으면 판례에 따름)
\begin{choices}
  \item 사단법인 이사의 대표권 제한은 이를 등기하지 않으면 악의의 제3자에게도 대항하지 못한다.
  \item 재단법인의 정관변경은 그 변경방법을 정관에서 정한 때에도 주무관청의 허가를 얻지 않으면 그 효력이 없다.
  \item 재단법인의 기본재산에 관한 근저당권 설정행위는 특별한 사정이 없는 한 주무관청의 허가를 얻을 필요가 없다.
  \item 재단법인의 기본재산 변경 시, 그로 인하여 기본재산이 새로이 편입되는 경우에는 주무관청의 허가를 얻을 필요가 없다.
  \item 법인에 대한 청산종결등기가 경료된 경우에도 청산사무가 종결되지 않는 한 그 범위 내에서는 청산법인으로서 존속한다.
\end{choices}
\end{question}
\answer{4}
